%! Direct, Inverse \& Joint Variation — Unit 4, Lesson 4.7 — Lesson Plan
\documentclass[10pt]{article}
\usepackage{algebra2-article}
\usepackage{algebra2-boxes}
\usepackage{graphicx}   % -article does NOT load graphicx; needed for thumbnails/screenshots
\graphicspath{{images/}}
\newcommand{\TallMath}[1]{$\displaystyle #1\rule[-1.4em]{0pt}{3.2em}$}
\newcommand{\CourseName}{Algebra 2}
\newcommand{\MeetingLength}{55 minutes}
\newcommand{\UnitNumberName}{Unit 4: Rational Functions \quad}
\newcommand{\LessonNumberName}{Lesson 4.7: Direct, Inverse \& Joint Variation}

\begin{document}
\begin{center}
    {\Large\bfseries \CourseName \\
    {\small\bfseries \UnitNumberName \LessonNumberName}}
\end{center}

\begin{tcolorbox}[colback=forestbg, colframe=forest, boxrule=0.9pt, arc=2mm,
    left=2mm, right=2mm, top=1.5mm, bottom=1.5mm]
    \textbf{Primary Objective:} Students can decide, from a table of values, whether two quantities are
    \textbf{directly proportional}, \textbf{inversely proportional}, or \textbf{neither}, by testing
    whether the quotients $y\div x$ or the products $x\cdot y$ are constant across \emph{every} row.
    They can find the \textbf{constant of variation} $k$, write the equation ($y=kx$ or $y=\frac kx$),
    and use it to answer the question asked; \emph{recognize} each graph --- a line through the
    \emph{origin} versus Lesson 4.4's hyperbola with the axes as asymptotes; and \emph{interpret} $k$, a
    point, and each asymptote in the language of the situation, including which perfectly valid
    solutions the \emph{context} rejects.
    \\[2pt]
    \textbf{Standards (2023 VA SOL):} \textbf{A2.F.1d} --- determine when two variables are directly
    proportional, inversely proportional, or neither, given a table of values; write an equation and
    create a graph to represent a direct or inverse variation, including situations in context.
    \textbf{A2.F.2f} --- for any $x$ in the domain of $f$, determine $f(x)$ using a graph or equation,
    and explain the meaning of $x$ and $f(x)$ in context. Applied: \textbf{A2.F.1a/c} (the rational
    parent and the vertical stretch that produces $y=\frac kx$, from 4.4) and \textbf{A2.F.2h}
    (the asymptotes, from 4.0). Builds on \textbf{A2.EI.4c} (5.6, interpreting and rejecting solutions).
    \emph{Joint and combined variation are not named in A2.F.1d --- they are carried here as clearly
    labeled enrichment and are not assessed.}
\end{tcolorbox}

\renewcommand{\arraystretch}{1.4}
\begin{skillbox}[Priority Ideas \& Skills]{goldbox}
\begin{tabularx}{\linewidth}{@{}X|X@{}}
\textbf{Priority Ideas \& Skills} & \textbf{Key Understandings (the ``why'')} \\[2pt]
\hline
\vspace{-6pt}
\begin{itemize}[leftmargin=1.1em, nosep, after=\vspace{2pt}]
  \item Run \emph{both} tests on a table: the quotients $y\div x$ and the products $x\cdot y$.
  \item Name the relationship --- direct, inverse, or neither --- from which row is constant.
  \item Find $k$ and write $y=kx$ or $y=\frac kx$; then answer the question that was asked.
  \item Use the shortcut when only the missing value is wanted: $\frac{y_1}{x_1}=\frac{y_2}{x_2}$ or $x_1y_1=x_2y_2$.
  \item Identify each graph on sight, and locate $k$ on it.
  \item State the equations of an inverse variation's asymptotes and say what each one means.
  \item Interpret $k$ with its units, and $f(x)$ at a point, in context.
  \item Reject the answers the \emph{situation} forbids --- and name that reason correctly.
\end{itemize}
&
\vspace{-6pt}
\begin{itemize}[leftmargin=1.1em, nosep, after=\vspace{2pt}]
  \item \textbf{One pair of numbers carries no information.} Every pair has a quotient and a product; only a table can decide.
  \item ``Goes up together'' is not direct and ``one up, one down'' is not inverse. A constant \emph{rate of change} makes a line; only a line through the origin is a variation.
  \item ``Neither'' is a correct, complete answer --- and out in the world it is the common one.
  \item An inverse variation \emph{is} the rational parent: $y=\frac kx$ is $\frac1x$ stretched by $k$, which is why $k$ cannot move the asymptotes.
  \item The asymptotes stop being features of a picture and start being sentences about a trip, a rectangle, a seesaw.
  \item $k$ always has units, and the units are how you know what $k$ means.
  \item 4.6 rejected answers because the algebra forbade them. Today the algebra is fine and the situation objects. Different reasons, different words.
\end{itemize}
\end{tabularx}
\end{skillbox}

\begin{skillbox}[{Vocabulary, Concepts, \& Theorems}]{sky}
\renewcommand{\arraystretch}{1.3}
\begin{tabularx}{\linewidth}{@{}l X@{}}
\textbf{Direct variation} & $y=kx$, $k\neq0$. The quotient $\frac yx$ is the same for every pair. Double $x$ and $y$ doubles. \\
\textbf{Inverse variation} & $y=\frac kx$, $k\neq0$. The product $xy$ is the same for every pair. Double $x$ and $y$ is halved. \\
\textbf{Constant of variation} & The number $k$ that stays fixed --- the constant quotient, or the constant product. In context it carries units and a meaning. \\
\textbf{Neither} & Every other relationship, including every line that misses the origin. A legitimate answer. \\
\textbf{Direct graph} & A line through the \emph{origin} with slope $k$. No excluded inputs, no asymptotes. \\
\textbf{Inverse graph} & The Lesson 4.4 parent stretched by $k$: \TallMath{y=\frac kx,\quad \text{VA } x=0,\quad \text{HA } y=0} \\
\textbf{Joint variation} \emph{(enrichment)} & $y=kxz$ --- direct in two variables at once. Combined variation, $y=\frac{kx}{z}$, is direct in one and inverse in another. \\
\end{tabularx}
\end{skillbox}

\begin{fixedskillbox}[Activate Prior Knowledge \& Spiral Review]{forestbg}
    The Warm-Up is the lesson in disguise, so debrief it before naming anything. \textbf{(1)} Two
    tables, and the only instruction is to fill in a \emph{quotient} row and a \emph{product} row. In
    Table P the quotients settle on $5$; in Table Q the products settle on $60$, and Q's quotient row is
    deliberately ugly ($15$, $3.75$, $1.\overline{6}$, $0.6$) so ``not constant'' is unmistakable. Get
    the two words on the board and leave them there --- do not say ``direct'' or ``inverse'' yet.
    \textbf{(2)} The rational parent $y=\frac6x$ from 4.4: three evaluations and both asymptotes. This
    is here because \emph{inverse variation is the parent function}, and when $t=\frac{240}{r}$ shows up
    in the notes the asymptotes must read as old news. \textbf{(3)} Two proportions solved by
    cross-multiplying, straight from 4.6 --- the fast route through every ``find the missing value''
    problem, because a direct variation is exactly the statement that two quotients are equal. Watch
    3(b), where the variable sits in the denominator.
\end{fixedskillbox}

\begin{skillbox}[Hook]{forestbg}
    Two tables on the board (they are in the notes). Table A: $x=1,2,3,4$ with $y=6,12,18,24$. Table B:
    $x=1,2,3,4$ with $y=9,12,15,18$. Point out what is true of \emph{both}: $y$ climbs by the same
    amount every step, so both are lines. \textbf{``Both of these are direct variations, right?''} Let
    them agree. Then run the doubling test out loud: take $x$ from $2$ to $4$. In A, $y$ goes
    $12\to24$ --- doubled. In B, $y$ goes $12\to18$ --- not even close. Fill in the quotient rows: A
    gives $6,6,6,6$ and B gives $9,6,5,4.5$. Finally ask for $y$ at $x=0$: A says $0$, B says $6$. That
    ``$+6$'' is the entire difference, and it is a \emph{fee you pay before you buy anything}. Post the
    conclusion and leave it up: \textbf{a direct variation is a line through the origin --- ``climbs
    steadily'' is not the test.} The Hook is built the way 4.6's was: identical surface features,
    different verdicts, so nobody can pattern-match their way past the actual criterion.
\end{skillbox}

\begin{skillbox}[Lesson]{forestbg}
\begin{multicols}{2}
\textbf{1. The two tests, and only the two tests.} Put the definitions up as arithmetic before they are
words: constant quotient $\Rightarrow$ direct, $y=kx$; constant product $\Rightarrow$ inverse,
$y=\frac kx$; neither constant $\Rightarrow$ neither. Say the sentence that prevents most of the
year's damage on this topic: \textbf{every pair of numbers has a quotient and a product, so one row can
never tell you which kind you have.} Then the two shapes side by side --- a line through the origin, and
the hyperbola. Students should be able to name the type from the picture alone.

\textbf{2. Direct, worked (the gasoline table).} Quotients all $3.40$, so $c=3.40g$. Do not stop at the
equation: ask for the \emph{units} of $k$ (dollars per gallon) and then for the sentence (``$k$ is the
price of one gallon''). Every direct variation passes through $(0,0)$, and here that is a claim about
the world --- zero gallons costs zero dollars. It is also the fastest disqualifier: a plan with a fee is
not a direct variation.

\textbf{3. Inverse, worked (the $240$-mile trip) --- and the unit closes its loop.} Products all $240$,
so $t=\frac{240}{r}$, and $k$ is not a rate this time; it is the \emph{distance}, which is just
$rt=d$ in disguise. Then the graph, and this is the part to protect: $t=\frac{240}{r}$ is the 4.4 parent
stretched by $240$. Both asymptotes now mean something sayable. ``$r=0$ is a vertical asymptote'' earns
nothing next to \emph{``standing still, you never get there''}; ``$t=0$ is a horizontal asymptote'' is
worth nothing next to \emph{``however fast you drive, the trip cannot take zero time.''} Close with the
domain in context: the equation allows $r=-30$, the road does not.

\columnbreak

\textbf{4. ``Neither'' is an answer (A2.F.1d says so).} Two traps, in opposite directions. The phone
plan climbs steadily and is not direct ($y=5x+15$). The second table \emph{falls} steadily and is not
inverse ($y=10-x$; the products are $9,16,21,24$). Expect a third of the room to call the second one
inverse. This section is the whole reason A2.F.1d includes the word ``neither,'' and students will not
volunteer it on their own.

\textbf{5. Find $k$, then answer.} The two-step routine, once each way, then the shortcut --- and name
it as yesterday's: a direct variation is literally $\frac{y_1}{x_1}=\frac{y_2}{x_2}$, a proportion, so
cross-multiply; an inverse variation is $x_1y_1=x_2y_2$ because both products equal $k$. Students shaky
on 4.6 get a second pass at cross-multiplying with friendly numbers. Insist they still write $k$ whenever
the problem asks what it \emph{means}.

\textbf{6. Joint and combined (enrichment --- label it out loud).} $y=kxz$ and $y=\frac{kx}{z}$: the
same move, done twice. A2.F.1d names direct and inverse only, so tell them plainly it will not be
assessed. Teach it if the room is with you; cut it without guilt if not. Nothing downstream depends on
it, and the triangle-area example ($k=\frac12$) pays for itself in thirty seconds.
\end{multicols}
\end{skillbox}

\begin{skillbox}[Explicit Instruction: Name It, Then Use It]{forestbg}
\begin{tabularx}{\linewidth}{@{}p{0.44\linewidth} X@{}}
\vspace{0pt}
\textbf{Post these. Both rows, every time.}
\begin{enumerate}[leftmargin=1.3em, nosep, after=\vspace{2pt}]
  \item \textbf{Test.} Compute $y\div x$ across the whole table. Compute $x\cdot y$ across the whole table.
  \item \textbf{Name it.} Quotients constant $\Rightarrow$ direct. Products constant $\Rightarrow$ inverse. Neither $\Rightarrow$ \emph{neither}.
  \item \textbf{Write it.} $y=kx$ \; or \; $y=\dfrac kx$, with the $k$ you just found.
  \item \textbf{Use it, then say it.} Substitute to answer the question --- and give $k$'s units and meaning.
\end{enumerate}

\vspace{3pt}
\textbf{The sentence:} \emph{one pair of numbers has both a quotient and a product, so one pair proves
nothing.}
&
\vspace{0pt}
\textbf{Direct.} \; $g:\;4,\,7,\,10,\,12$ \quad $c:\;13.60,\,23.80,\,34.00,\,40.80$
\[ \frac cg = 3.40 \text{ every time} \;\Rightarrow\; c=3.40g \]
\centerline{\small $k=\$3.40$ \emph{per gallon}; \; the graph runs through $(0,0)$.}

\vspace{4pt}
\textbf{Inverse.} \; $r:\;30,\,40,\,48,\,60$ \quad $t:\;8,\,6,\,5,\,4$
\[ rt = 240 \text{ every time} \;\Rightarrow\; t=\frac{240}{r} \]
\centerline{\small $k=240$ \emph{miles}; \; VA $r=0$, \; HA $t=0$.}

\vspace{4pt}
\textbf{The two questions students must answer aloud:} \emph{which row came out constant?} and
\emph{what does $k$ mean here, with units?}
\end{tabularx}
\end{skillbox}

\begin{skillbox}[Active Monitoring]{redbox}
Circulate during the guided practice and Tier work. Watch for and correct:
\begin{itemize}[leftmargin=1.2em, nosep]
  \item \textbf{$k$ computed from one row} and never tested against the rest of the table --- the signature error of this lesson, and the whole basis of activity Tier A part 3;
  \item \textbf{only the row that matches the guess} filled in, which is guessing dressed up as testing;
  \item \textbf{``$y$ goes down, so it's inverse''} --- decreasing is not the test (notes \S4, exit ticket 2);
  \item \textbf{``it goes up by the same amount, so it's direct''} --- that makes it a \emph{line}, not a variation (the Hook, activity Tier R part 3);
  \item \textbf{$k$ found correctly, then written into the wrong form} --- $y=kx$ for an inverse relationship (homework 5);
  \item \textbf{``neither'' refused}, and the table re-checked three times in search of a $k$;
  \item \textbf{$k$ reported as a bare number} with no units and no meaning --- half the credit on every context item;
  \item \textbf{a negative $k$ treated as an error} (homework 2(f)) --- the line simply falls;
  \item \textbf{``extraneous'' used for a context rejection} --- the 4.6 word, misapplied; today the algebra has no objection at all.
\end{itemize}
Cold-call prompts: ``Which row did you test --- both of them?'' ``Show me the product for the
\emph{last} pair, not the first.'' ``What are the units of your $k$?'' ``Say what $k$ means in a
sentence with no letters in it.'' ``Where is $k$ on that graph?'' ``What does the vertical asymptote
say about the actual situation?'' ``You threw that answer out --- was it the algebra or the situation
that made you?''
\end{skillbox}

\begin{skillbox}[Group Work \& Differentiation]{redbox}
\textbf{Constant Quotient, Constant Product} activity (tiered; mirrors the notes).
\begin{multicols}{3}
\textbf{Tier R --- Remediate.} Volume, then contrast. \emph{Part 1:} four tables, both arithmetic rows
required on each --- (a) direct $k=4$, (b) inverse $k=48$, (c) \textbf{neither} (decreasing, and the
products $20,34,42,44$ nearly stabilize, which is exactly why ``every row'' matters), (d) direct $k=7$.
\emph{Part 2:} four find-$k$-then-answer items, two each way. \emph{Part 3:} \textbf{twins} --- two
tables that both climb $5$ for every $2$, one direct ($y=2.5x$) and one not ($y=2.5x+3$). That pair is
the tier's argument: nothing in the \emph{pattern} predicted which. If a group finishes early, ask them
to predict $y$ at $x=0$ for both before you say a word.

\columnbreak
\textbf{Tier A --- Approaching.} \emph{Part 1:} a spring, $s=0.75w$ --- the units of $k$ (cm per pound)
and why the graph must pass through $(0,0)$. \emph{Part 2:} the \textbf{A2.F.2f graph item} --- all
rectangles of area $36$, with $h=\frac{36}{b}$ pre-drawn. Read $(6,6)$ and recognize the square; find
$b$ when $h=8$; say what the flattening tail means (as short as you like, never flat); and reject
$b=-4$ \emph{for the situation}, not as extraneous. \emph{Part 3:} an \textbf{error analysis where the
arithmetic is flawless} --- the student computes $k=\frac62=3$ from the first pair, writes $y=3x$, and
never tests row two ($y=2x+2$). Expect a group to defend them; that is the right conversation.

\columnbreak
\textbf{Tier E --- Extension.} \emph{Part 1:} joint and combined variation, including deriving
$V=\pi r^2h$ from ``$V$ varies jointly as $h$ and $r^2$'' --- they discover the formula rather than
receive it. \emph{Part 2:} \textbf{the lesson with no numbers in it} --- a scaling table (triple $x$,
halve $x$, double $x$ and triple $z$\dots) closed by an actual proof that replacing $x$ with $3x$ in
$y=\frac kx$ gives $\frac13 y$. \emph{Part 3:} the deepest idea available here --- $y=\frac6x$ is an
\emph{inverse} variation in $x$ and a \emph{direct} variation in $\frac1x$, which is precisely why
$y=\frac kx$ is the parent $\frac1x$ scaled by $k$; then a design task (an inverse and a direct through
$(4,5)$) and the impossible one (a direct variation through $(0,3)$).
\end{multicols}
\end{skillbox}

\begin{skillbox}[Individual Work \& Assessment]{redbox}
\textbf{Exit Ticket} (independent, no notes): \textbf{(1)} \emph{name it and use it} on the table
$x=2,3,4,6$ / $y=18,12,9,6$ --- both arithmetic rows are printed, so a student who fills in only the
product row has guessed; products all $36$, so $y=\frac{36}{x}$, giving $y=4.5$ at $x=8$ and $x=12$ at
$y=3$, with asymptotes $x=0$ and $y=0$ (the unit's closing loop: the 4.4 parent stretched by $36$, and
$k$ cannot move an asymptote). \textbf{(2)} the \textbf{written-reasoning item} --- a classmate argues
``$y$ goes down when $x$ goes up, so it is inverse,'' on a table whose products are $12,18,18,12$;
\emph{two of them match on purpose}, so a student who checks one or two pairs and stops will agree with
the classmate. Grade the \emph{reason}: full credit names the constant-product test and reports that it
fails across all four pairs. \textbf{(3)} an \textbf{SOL-style multiple-choice} item --- a pump filling
a pool, $k=rt=8\cdot45=360$ gallons, so a $12$ gal/min pump takes $\frac{360}{12}=30$ minutes
(\textbf{A}); \textbf{B} $=67.5$ sets up a direct proportion, \textbf{C} $=3.75$ divides $45$ by the new
rate without ever finding $k$, and \textbf{D} $=22.5$ halves $45$ on instinct. The sense check students
should own by now --- a faster pump takes \emph{less} time --- kills B on sight. Sort into ``got it /
named it from one row / used the wrong test / computed it but cannot say what $k$ means.'' The last pile
is the one to fix before the unit test, where interpretation carries as many points as the arithmetic.
\end{skillbox}

\begin{skillbox}[Reinforcement \& Extension]{goldbox}
\textbf{Homework:} four \textbf{tables to name} (direct $2.5$, inverse $24$, \textbf{neither}
($y=3x+1$), direct $7$); six \textbf{find-$k$-and-answer} items running both directions, closing with a
\textbf{negative $k$} so that ``$k<0$'' never reads as ``I did it wrong''; three \textbf{decide from the
sentence} items, including a taxi with a $\$3.00$ flag drop --- the same $+b$ as problem 1(c), now in
English; \textbf{two graphs to read} (A2.F.2f), one line through the origin and one hyperbola, ending on
the structural question \emph{why can a direct variation never have an asymptote?} (it is a polynomial:
no excluded input to build a wall at, no value it settles toward --- the 4.0 compare-and-contrast table,
and fair game on the test); an \textbf{error analysis} on the lesson's other big mistake, a correct $k$
poured into the wrong \emph{form} ($y=24x$ instead of $y=\frac{24}{x}$, giving $288$ when the sanity
check says the answer must be smaller than $8$); and \textbf{the seesaw} ($d=\frac{300}{w}$, $k$ in
pound-feet) as the capstone --- both asymptotes interpreted in English, an $8$-foot beam turning into
the extra restriction $\frac{300}{w}\le8 \Rightarrow w\ge37.5$, and finally $w=-60$, which \emph{solves
the equation} and is thrown out anyway. That last part is the reason this lesson follows 4.6, and
``it's extraneous'' is not an acceptable answer to it. \textbf{Extension:} joint and combined variation;
two scaling claims to \emph{prove} rather than assert; and the sharpest item in the packet --- could a
table be direct \emph{and} inverse at once? ($kx=\frac mx$ forces $x^2=\frac mk$, so at most two
$x$-values; a four-row table cannot.) It lands on a direct and an inverse variation through $(2,6)$
meeting again at $(-2,-6)$. \textbf{Preview:} the \textbf{Unit 4 Test} --- this is the last lesson of the
unit. The packet's practice test mirrors the real one in format and ideas with different numbers; the
habit worth more than any formula on it is \emph{factor first, restrict early, and check what the answer
is supposed to mean.}
\end{skillbox}
\end{document}
